\documentclass[10pt,a4paper]{amsart}
\usepackage{amsmath,amssymb,amsthm,graphicx}
\usepackage{verbatim}
\newenvironment{reference}{\comment}{\endcomment}
%\newenvironment{reference}{}{}
\newenvironment{slogan}{\comment}{\endcomment}
\newenvironment{history}{\comment}{\endcomment}

% For commutative diagrams we use Xy-pic
\usepackage[all]{xy}

% We use 2cell for 2-commutative diagrams.
\xyoption{2cell}
% Theorem environments.
%
\theoremstyle{plain}
\newtheorem{theorem}[subsection]{Theorem}
\newtheorem{proposition}[subsection]{Proposition}
\newtheorem{lemma}[subsection]{Lemma}

\theoremstyle{definition}
\newtheorem{definition}[subsection]{Definition}
\newtheorem{example}[subsection]{Example}
\newtheorem{exercise}[subsection]{Exercise}
\newtheorem{situation}[subsection]{Situation}

\theoremstyle{remark}
\newtheorem{remark}[subsection]{Remark}
\newtheorem{remarks}[subsection]{Remarks}

\numberwithin{equation}{subsection}

% Macros
%
\def\lim{\mathop{\mathrm{lim}}\nolimits}
\def\colim{\mathop{\mathrm{colim}}\nolimits}
\def\Spec{\mathop{\mathrm{Spec}}}
\def\Hom{\mathop{\mathrm{Hom}}\nolimits}
\def\Ext{\mathop{\mathrm{Ext}}\nolimits}
\def\SheafHom{\mathop{\mathcal{H}\!\mathit{om}}\nolimits}
\def\SheafExt{\mathop{\mathcal{E}\!\mathit{xt}}\nolimits}
\def\Sch{\mathit{Sch}}
\def\Mor{\mathop{\mathrm{Mor}}\nolimits}
\def\Ob{\mathop{\mathrm{Ob}}\nolimits}
\def\Sh{\mathop{\mathit{Sh}}\nolimits}
\def\NL{\mathop{N\!L}\nolimits}
\def\CH{\mathop{\mathrm{CH}}\nolimits}
\def\proetale{{pro\text{-}\acute{e}tale}}
\def\etale{{\acute{e}tale}}
\def\QCoh{\mathit{QCoh}}
\def\Ker{\mathop{\mathrm{Ker}}}
\def\Im{\mathop{\mathrm{Im}}}
\def\Coker{\mathop{\mathrm{Coker}}}
\def\Coim{\mathop{\mathrm{Coim}}}

% Boxtimes
%
\DeclareMathSymbol{\boxtimes}{\mathbin}{AMSa}{"02}

%
% Macros for moduli stacks/spaces
%
\def\QCohstack{\mathcal{QC}\!\mathit{oh}}
\def\Cohstack{\mathcal{C}\!\mathit{oh}}
\def\Spacesstack{\mathcal{S}\!\mathit{paces}}
\def\Quotfunctor{\mathrm{Quot}}
\def\Hilbfunctor{\mathrm{Hilb}}
\def\Curvesstack{\mathcal{C}\!\mathit{urves}}
\def\Polarizedstack{\mathcal{P}\!\mathit{olarized}}
\def\Complexesstack{\mathcal{C}\!\mathit{omplexes}}
% \Pic is the operator that assigns to X its picard group, usage \Pic(X)
% \Picardstack_{X/B} denotes the Picard stack of X over B
% \Picardfunctor_{X/B} denotes the Picard functor of X over B
\def\Pic{\mathop{\mathrm{Pic}}\nolimits}
\def\Picardstack{\mathcal{P}\!\mathit{ic}}
\def\Picardfunctor{\mathrm{Pic}}
\def\Deformationcategory{\mathcal{D}\!\mathit{ef}}
\usepackage{cleveref}
\crefname{originallemma}{original lemma}{original lemmas}
\Crefname{originallemma}{Original lemma}{Original lemmas}
\makeatletter
\newlabel{coherent-lemma-characterize-finite}{{02OG}{source}{}{sourcepointer-1}{}}
\newlabel{coherent-lemma-characterize-finite@cref}{{[originallemma][0][]02OG}{[1][0][]source}{}{}{}}
\makeatother
\crefname{originalreference}{original reference}{original references}
\Crefname{originalreference}{Original reference}{Original references}
\makeatletter
\newlabel{example-rational-bridge}{{0E3M}{source}{}{sourcepointer-2}{}}
\newlabel{example-rational-bridge@cref}{{[originalreference][0][]0E3M}{[1][0][]source}{}{}{}}
\makeatother
\crefname{originalreference}{original reference}{original references}
\Crefname{originalreference}{Original reference}{Original references}
\makeatletter
\newlabel{example-rational-tail}{{0E3H}{source}{}{sourcepointer-3}{}}
\newlabel{example-rational-tail@cref}{{[originalreference][0][]0E3H}{[1][0][]source}{}{}{}}
\makeatother
\crefname{originalreference}{original reference}{original references}
\Crefname{originalreference}{Original reference}{Original references}
\makeatletter
\newlabel{lemma-automatic}{{0BY5}{source}{}{sourcepointer-4}{}}
\newlabel{lemma-automatic@cref}{{[originalreference][0][]0BY5}{[1][0][]source}{}{}{}}
\makeatother
\crefname{originalreference}{original reference}{original references}
\Crefname{originalreference}{Original reference}{Original references}
\makeatletter
\newlabel{lemma-closed-subscheme-nodal-curve}{{0E38}{source}{}{sourcepointer-5}{}}
\newlabel{lemma-closed-subscheme-nodal-curve@cref}{{[originalreference][0][]0E38}{[1][0][]source}{}{}{}}
\makeatother
\crefname{originalreference}{original reference}{original references}
\Crefname{originalreference}{Original reference}{Original references}
\makeatletter
\newlabel{lemma-contract-gorenstein-canonical}{{0E7P}{source}{}{sourcepointer-6}{}}
\newlabel{lemma-contract-gorenstein-canonical@cref}{{[originalreference][0][]0E7P}{[1][0][]source}{}{}{}}
\makeatother
\crefname{originalreference}{original reference}{original references}
\Crefname{originalreference}{Original reference}{Original references}
\makeatletter
\newlabel{lemma-contracting-rational-bridges}{{0E3Q}{source}{}{sourcepointer-7}{}}
\newlabel{lemma-contracting-rational-bridges@cref}{{[originalreference][0][]0E3Q}{[1][0][]source}{}{}{}}
\makeatother
\crefname{originalreference}{original reference}{original references}
\Crefname{originalreference}{Original reference}{Original references}
\makeatletter
\newlabel{lemma-contracting-rational-tails}{{0E3L}{source}{}{sourcepointer-8}{}}
\newlabel{lemma-contracting-rational-tails@cref}{{[originalreference][0][]0E3L}{[1][0][]source}{}{}{}}
\makeatother
\crefname{originalreference}{original reference}{original references}
\Crefname{originalreference}{Original reference}{Original references}
\makeatletter
\newlabel{lemma-no-rational-tail-semiample-genus-geq-2}{{0E3K}{source}{}{sourcepointer-9}{}}
\newlabel{lemma-no-rational-tail-semiample-genus-geq-2@cref}{{[originalreference][0][]0E3K}{[1][0][]source}{}{}{}}
\makeatother
\crefname{originalreference}{original reference}{original references}
\Crefname{originalreference}{Original reference}{Original references}
\makeatletter
\newlabel{lemma-rational-bridge-canonical}{{0E3N}{source}{}{sourcepointer-10}{}}
\newlabel{lemma-rational-bridge-canonical@cref}{{[originalreference][0][]0E3N}{[1][0][]source}{}{}{}}
\makeatother
\crefname{originalreference}{original reference}{original references}
\Crefname{originalreference}{Original reference}{Original references}
\makeatletter
\newlabel{lemma-rational-tail-negative}{{0E63}{source}{}{sourcepointer-11}{}}
\newlabel{lemma-rational-tail-negative@cref}{{[originalreference][0][]0E63}{[1][0][]source}{}{}{}}
\makeatother
\crefname{originalreference}{original reference}{original references}
\Crefname{originalreference}{Original reference}{Original references}
\makeatletter
\newlabel{lemma-vanishing-twist}{{0B62}{source}{}{sourcepointer-12}{}}
\newlabel{lemma-vanishing-twist@cref}{{[originalreference][0][]0B62}{[1][0][]source}{}{}{}}
\makeatother
\crefname{originallemma}{original lemma}{original lemmas}
\Crefname{originallemma}{Original lemma}{Original lemmas}
\makeatletter
\newlabel{more-morphisms-lemma-check-h1-fibre-zero}{{0E7F}{source}{}{sourcepointer-13}{}}
\newlabel{more-morphisms-lemma-check-h1-fibre-zero@cref}{{[originallemma][0][]0E7F}{[1][0][]source}{}{}{}}
\makeatother
\crefname{originallemma}{original lemma}{original lemmas}
\Crefname{originallemma}{Original lemma}{Original lemmas}
\makeatletter
\newlabel{more-morphisms-lemma-h1-fibre-zero-check-h0-kappa}{{0E7I}{source}{}{sourcepointer-14}{}}
\newlabel{more-morphisms-lemma-h1-fibre-zero-check-h0-kappa@cref}{{[originallemma][0][]0E7I}{[1][0][]source}{}{}{}}
\makeatother
\crefname{originaltheorem}{original theorem}{original theorems}
\Crefname{originaltheorem}{Original theorem}{Original theorems}
\makeatletter
\newlabel{more-morphisms-theorem-stein-factorization-Noetherian}{{03H0}{source}{}{sourcepointer-15}{}}
\newlabel{more-morphisms-theorem-stein-factorization-Noetherian@cref}{{[originaltheorem][0][]03H0}{[1][0][]source}{}{}{}}
\makeatother
\crefname{originallemma}{original lemma}{original lemmas}
\Crefname{originallemma}{Original lemma}{Original lemmas}
\makeatletter
\newlabel{morphisms-lemma-proper-ample-is-proj}{{0C6J}{source}{}{sourcepointer-16}{}}
\newlabel{morphisms-lemma-proper-ample-is-proj@cref}{{[originallemma][0][]0C6J}{[1][0][]source}{}{}{}}
\makeatother
\crefname{originallemma}{original lemma}{original lemmas}
\Crefname{originallemma}{Original lemma}{Original lemmas}
\makeatletter
\newlabel{properties-lemma-criterion-reduced}{{0344}{source}{}{sourcepointer-17}{}}
\newlabel{properties-lemma-criterion-reduced@cref}{{[originallemma][0][]0344}{[1][0][]source}{}{}{}}
\makeatother
\crefname{originalsection}{original section}{original sections}
\Crefname{originalsection}{Original section}{Original sections}
\makeatletter
\newlabel{section-curves-function-fields}{{0BXX}{source}{}{sourcepointer-18}{}}
\newlabel{section-curves-function-fields@cref}{{[originalsection][0][]0BXX}{[1][0][]source}{}{}{}}
\makeatother
\title{Characterization of the stable contraction of a nodal curve of genus at least 2 with H\textasciicircum{}0(X, O\_X) = k}
\author{The Stacks Project Authors}
\date{stacks-project:curves.tex:local-snapshot-20261012; GNU FDL 1.2 or later}
\begin{document}
\maketitle
\noindent\textit{Editorial note.} All delivered mathematical text is the author's own, copied from the pinned Stacks Project chapter file curves.tex; no statement or proof step was rewritten, reordered or removed.\par\medskip
\noindent\textit{Editorial note.} The counted proof is the author's complete argument: existence by the composition of tail and bridge contractions, the full fibre analysis of the competitor, the rational-tail-to-point claim with both cases, the induction over tail contractions, and the Proj comparison closing the uniqueness proof, including the author's remark that the uniqueness of the isomorphism is omitted as a verification.\par\medskip
\noindent\textit{Editorial note.} The label lemma-characterize-contraction-to-stable names the counted statement, and the proof environment immediately following it in the source is delivered verbatim as its proof.\par\medskip
\noindent\textit{Editorial note.} The counted statement is carried by the author's own lemma environment in the chapter source.\par\medskip
\noindent\textit{Allowed standard prerequisites (editorial).}\begin{itemize}
\item Contractions of rational tails and rational bridges. Same chapter, Lemmas lemma-contracting-rational-tails and lemma-contracting-rational-bridges with Examples example-rational-tail and example-rational-bridge. Supplies the existence half: a morphism c with O\_Y = c\_*O\_X, R\textasciicircum{}1c\_*O\_X = 0 and ample omega\_Y, obtained as a composition of tail and bridge contractions.
\item Stein factorization and fibres of c. More on Morphisms, Theorem more-morphisms-theorem-stein-factorization-Noetherian and Lemmas more-morphisms-lemma-check-h1-fibre-zero, more-morphisms-lemma-h1-fibre-zero-check-h0-kappa. Gives geometrically connected fibres and the identities H\textasciicircum{}1(X\_y, O) = 0, H\textasciicircum{}0(X\_y, O) = kappa(y) used to classify the fibres of the competing contraction.
\end{itemize}\medskip
\ifdefined\hypertarget\hypertarget{sourcepointer-1}{}\fi
\noindent\textit{Source reference.} Original lemma 02OG: sources/stacks-chow/tags.txt, source lines 3515--3515. Prerequisite pointer invoked at proof line(s) 6721; it names a standard fact of a sibling chapter that the counted argument uses, and every step of the counted proof built on it is retained verbatim in the delivered body.\par\medskip
\ifdefined\hypertarget\hypertarget{sourcepointer-2}{}\fi
\noindent\textit{Source reference.} Original reference 0E3M: sources/stacks-curves/pinned-original-source.tex, source lines 6201--6201. Prerequisite pointer invoked at proof line(s) 6681; it names a standard fact of the same chapter that the counted argument uses, and every step of the counted proof built on it is retained verbatim in the delivered body.\par\medskip
\ifdefined\hypertarget\hypertarget{sourcepointer-3}{}\fi
\noindent\textit{Source reference.} Original reference 0E3H: sources/stacks-curves/pinned-original-source.tex, source lines 5933--5933. Prerequisite pointer invoked at proof line(s) 6680; it names a standard fact of the same chapter that the counted argument uses, and every step of the counted proof built on it is retained verbatim in the delivered body.\par\medskip
\ifdefined\hypertarget\hypertarget{sourcepointer-4}{}\fi
\noindent\textit{Source reference.} Original reference 0BY5: sources/stacks-curves/pinned-original-source.tex, source lines 1277--1277. Prerequisite pointer invoked at proof line(s) 6708; it names a standard fact of the same chapter that the counted argument uses, and every step of the counted proof built on it is retained verbatim in the delivered body.\par\medskip
\ifdefined\hypertarget\hypertarget{sourcepointer-5}{}\fi
\noindent\textit{Source reference.} Original reference 0E38: sources/stacks-curves/pinned-original-source.tex, source lines 4794--4794. Prerequisite pointer invoked at proof line(s) 6713; it names a standard fact of the same chapter that the counted argument uses, and every step of the counted proof built on it is retained verbatim in the delivered body.\par\medskip
\ifdefined\hypertarget\hypertarget{sourcepointer-6}{}\fi
\noindent\textit{Source reference.} Original reference 0E7P: sources/stacks-curves/pinned-original-source.tex, source lines 6572--6572. Prerequisite pointer invoked at proof line(s) 6771; it names a standard fact of the same chapter that the counted argument uses, and every step of the counted proof built on it is retained verbatim in the delivered body.\par\medskip
\ifdefined\hypertarget\hypertarget{sourcepointer-7}{}\fi
\noindent\textit{Source reference.} Original reference 0E3Q: sources/stacks-curves/pinned-original-source.tex, source lines 6491--6491. Prerequisite pointer invoked at proof line(s) 6671; it names a standard fact of the same chapter that the counted argument uses, and every step of the counted proof built on it is retained verbatim in the delivered body.\par\medskip
\ifdefined\hypertarget\hypertarget{sourcepointer-8}{}\fi
\noindent\textit{Source reference.} Original reference 0E3L: sources/stacks-curves/pinned-original-source.tex, source lines 6135--6135. Prerequisite pointer invoked at proof line(s) 6670; it names a standard fact of the same chapter that the counted argument uses, and every step of the counted proof built on it is retained verbatim in the delivered body.\par\medskip
\ifdefined\hypertarget\hypertarget{sourcepointer-9}{}\fi
\noindent\textit{Source reference.} Original reference 0E3K: sources/stacks-curves/pinned-original-source.tex, source lines 6080--6080. Prerequisite pointer invoked at proof line(s) 6768; it names a standard fact of the same chapter that the counted argument uses, and every step of the counted proof built on it is retained verbatim in the delivered body.\par\medskip
\ifdefined\hypertarget\hypertarget{sourcepointer-10}{}\fi
\noindent\textit{Source reference.} Original reference 0E3N: sources/stacks-curves/pinned-original-source.tex, source lines 6367--6367. Prerequisite pointer invoked at proof line(s) 6774; it names a standard fact of the same chapter that the counted argument uses, and every step of the counted proof built on it is retained verbatim in the delivered body.\par\medskip
\ifdefined\hypertarget\hypertarget{sourcepointer-11}{}\fi
\noindent\textit{Source reference.} Original reference 0E63: sources/stacks-curves/pinned-original-source.tex, source lines 5987--5987. Prerequisite pointer invoked at proof line(s) 6738; it names a standard fact of the same chapter that the counted argument uses, and every step of the counted proof built on it is retained verbatim in the delivered body.\par\medskip
\ifdefined\hypertarget\hypertarget{sourcepointer-12}{}\fi
\noindent\textit{Source reference.} Original reference 0B62: sources/stacks-curves/pinned-original-source.tex, source lines 1354--1354. Prerequisite pointer invoked at proof line(s) 6770; it names a standard fact of the same chapter that the counted argument uses, and every step of the counted proof built on it is retained verbatim in the delivered body.\par\medskip
\ifdefined\hypertarget\hypertarget{sourcepointer-13}{}\fi
\noindent\textit{Source reference.} Original lemma 0E7F: sources/stacks-chow/tags.txt, source lines 17780--17780. Prerequisite pointer invoked at proof line(s) 6702; it names a standard fact of a sibling chapter that the counted argument uses, and every step of the counted proof built on it is retained verbatim in the delivered body.\par\medskip
\ifdefined\hypertarget\hypertarget{sourcepointer-14}{}\fi
\noindent\textit{Source reference.} Original lemma 0E7I: sources/stacks-chow/tags.txt, source lines 17783--17783. Prerequisite pointer invoked at proof line(s) 6704; it names a standard fact of a sibling chapter that the counted argument uses, and every step of the counted proof built on it is retained verbatim in the delivered body.\par\medskip
\ifdefined\hypertarget\hypertarget{sourcepointer-15}{}\fi
\noindent\textit{Source reference.} Original theorem 03H0: sources/stacks-chow/tags.txt, source lines 4547--4547. Prerequisite pointer invoked at proof line(s) 6693; it names a standard fact of a sibling chapter that the counted argument uses, and every step of the counted proof built on it is retained verbatim in the delivered body.\par\medskip
\ifdefined\hypertarget\hypertarget{sourcepointer-16}{}\fi
\noindent\textit{Source reference.} Original lemma 0C6J: sources/stacks-chow/tags.txt, source lines 15299--15299. Prerequisite pointer invoked at proof line(s) 6784; it names a standard fact of a sibling chapter that the counted argument uses, and every step of the counted proof built on it is retained verbatim in the delivered body.\par\medskip
\ifdefined\hypertarget\hypertarget{sourcepointer-17}{}\fi
\noindent\textit{Source reference.} Original lemma 0344: sources/stacks-chow/tags.txt, source lines 4079--4079. Prerequisite pointer invoked at proof line(s) 6711; it names a standard fact of a sibling chapter that the counted argument uses, and every step of the counted proof built on it is retained verbatim in the delivered body.\par\medskip
\ifdefined\hypertarget\hypertarget{sourcepointer-18}{}\fi
\noindent\textit{Source reference.} Original section 0BXX: sources/stacks-curves/pinned-original-source.tex, source lines 264--264. Prerequisite pointer invoked at proof line(s) 6751; it names a standard fact of the same chapter that the counted argument uses, and every step of the counted proof built on it is retained verbatim in the delivered body.\par\medskip
\begin{lemma}
\label{lemma-characterize-contraction-to-stable}
Let $k$ be a field. Let $X$ be a proper scheme over $k$ of dimension
$1$ with $H^0(X, \mathcal{O}_X) = k$ having genus $g \geq 2$.
Assume the singularities of $X$ are at-worst-nodal.
There is a unique morphism (up to unique isomorphism)
$$
c : X \longrightarrow Y
$$
of schemes over $k$ having the following properties:
\begin{enumerate}
\item $Y$ is proper over $k$, $\dim(Y) = 1$, the singularities of $Y$
are at-worst-nodal,
\item $\mathcal{O}_Y = c_*\mathcal{O}_X$ and $R^1c_*\mathcal{O}_X = 0$, and
\item $\omega_Y$ is ample on $Y$.
\end{enumerate}
\end{lemma}

\begin{proof}
Existence: A morphism with all the properties listed exists by
combining Lemmas \ref{lemma-contracting-rational-tails} and
\ref{lemma-contracting-rational-bridges} as discussed in the
introduction to this section.
Moreover, we see that it can be written as a composition
$$
X \to X_1 \to X_2 \ldots \to X_n \to X_{n + 1} \to \ldots \to X_{n + n'}
$$
where the first $n$ morphisms are contractions of rational tails
and the last $n'$ morphisms are contractions of rational bridges.
Note that property (2) holds for each contraction of a rational
tail (Example \ref{example-rational-tail}) and contraction of a
rational bridge (Example \ref{example-rational-bridge}).
It is easy to see that this property is inherited by compositions of morphisms.

\medskip\noindent
Uniqueness: Let $c : X \to Y$ be a morphism satisfying conditions
(1), (2), and (3). We will show that there is a unique isomorphism
$X_{n + n'} \to Y$ compatible with the morphisms $X \to X_{n + n'}$ and $c$.

\medskip\noindent
Before we start the proof we make some observations about $c$.
We first observe that the fibres of $c$ are geometrically connected by
More on Morphisms, Theorem
\ref{more-morphisms-theorem-stein-factorization-Noetherian}.
In particular $c$ is surjective.
For a closed point $y \in Y$ the fibre $X_y$ satisfies
$$
H^1(X_y, \mathcal{O}_{X_y}) = 0
\quad\text{and}\quad
H^0(X_y, \mathcal{O}_{X_y}) = \kappa(y)
$$
The first equality by More on Morphisms, Lemma
\ref{more-morphisms-lemma-check-h1-fibre-zero}
and the second by More on Morphisms, Lemma
\ref{more-morphisms-lemma-h1-fibre-zero-check-h0-kappa}.
Thus either $X_y = x$ where $x$ is the unique point of $X$ mapping to
$y$ and has the same residue field as $y$, or $X_y$ is a $1$-dimensional
proper scheme over $\kappa(y)$. Observe that in the second case
$X_y$ is Cohen-Macaulay (Lemma \ref{lemma-automatic}).
However, since $X$ is reduced, we see that $X_y$ must be reduced
at all of its generic points (details omitted), and hence $X_y$
is reduced by Properties, Lemma \ref{properties-lemma-criterion-reduced}.
It follows that the singularities of $X_y$ are at-worst-nodal
(Lemma \ref{lemma-closed-subscheme-nodal-curve}).
Note that the genus of $X_y$ is zero (see above).
Finally, there are only a finite number of points $y$ where
the fibre $X_y$ has dimension $1$, say
$\{y_1, \ldots, y_r\}$, and $c^{-1}(Y \setminus \{y_1, \ldots, y_r\})$
maps isomorphically to $Y \setminus \{y_1, \ldots, y_r\}$ by $c$.
Some details omitted; hint: outside of $\{y_1, \ldots, y_r\}$
the morphism $c$ is finite, see
Cohomology of Schemes, Lemma \ref{coherent-lemma-characterize-finite}.

\medskip\noindent
Let $C \subset X$ be a rational tail.
We claim that $c$ maps $C$ to a point. Assume that this
is not the case to get a contradiction. Then the image
of $C$ is an irreducible component $D \subset Y$.
Recall that $H^0(C, \mathcal{O}_C) = k'$ is a finite separable
extension of $k$ and that $C$ has a $k'$-rational point $x$
which is also the unique intersection of $C$ with the ``rest'' of $X$.
We conclude from the general discussion above that
$C \setminus \{x\} \subset c^{-1}(Y \setminus \{y_1, \ldots, y_r\})$
maps isomorphically to an open $V$ of $D$. Let $y = c(x) \in D$.
Observe that $y$ is the only point of $D$ meeting the
``rest'' of $Y$. If $y \not \in \{y_1, \ldots, y_r\}$, then $C \cong D$
and it is clear that $D$ is a rational tail of $Y$ which is
a contradiction with the ampleness of $\omega_Y$
(Lemma \ref{lemma-rational-tail-negative}).
Thus $y \in \{y_1, \ldots, y_r\}$ and $\dim(X_y) = 1$.
Then $x \in X_y \cap C$ and $x$ is a smooth point of $X_y$ and $C$
(Lemma \ref{lemma-closed-subscheme-nodal-curve}).
If $y \in D$ is a singular point of $D$, then $y$ is a node
and then $Y = D$ (because there cannot be another component of
$Y$ passing through $y$ by Lemma \ref{lemma-closed-subscheme-nodal-curve}).
Then $X = X_y \cup C$ which means $g = 0$ because it is
equal to the genus of $X_y$ by the discussion in
Example \ref{example-rational-tail}; a contradiction.
If $y \in D$ is a smooth point of $D$, then
$C \to D$ is an isomorphism (because the nonsingular projective
model is unique and $C$ and $D$ are birational, see
Section \ref{section-curves-function-fields}). Then $D$ is
a rational tail of $Y$ which is a contradiction with
ampleness of $\omega_Y$.

\medskip\noindent
Assume $n \geq 1$. If $C \subset X$ is the rational tail contracted
by $X \to X_1$, then we see that $C$ is mapped to a point of $Y$ by
the previous paragraph. Hence $c : X \to Y$ factors through $X \to X_1$
(because $X$ is the pushout of $C$ and $X_1$, see discussion in
Example \ref{example-rational-tail}).
After replacing $X$ by $X_1$ we have decreased
$n$. By induction we may assume $n = 0$, i.e., $X$ does not have
a rational tail.

\medskip\noindent
Assume $n = 0$, i.e., $X$ does not have any rational tails.
Then $\omega_X^{\otimes 2}$ and $\omega_X^{\otimes 3}$ are
globally generated by Lemma \ref{lemma-no-rational-tail-semiample-genus-geq-2}.
It follows that $H^1(X, \omega_X^{\otimes 3}) = 0$ by
Lemma \ref{lemma-vanishing-twist}.
By Lemma \ref{lemma-contract-gorenstein-canonical} applied with $m = 3$
we find that $c^*\omega_Y \cong \omega_X$.
We also have that $\omega_X = (X \to X_{n'})^*\omega_{X_{n'}}$ by
Lemma \ref{lemma-rational-bridge-canonical} and induction.
Applying the projection formula for both $c$ and
$X \to X_{n'}$ we conclude that
$$
\Gamma(X_{n'}, \omega_{X_{n'}}^{\otimes m}) =
\Gamma(X, \omega_X^{\otimes m}) =
\Gamma(Y, \omega_Y^{\otimes m})
$$
for all $m$.
Since $X_{n'}$ and $Y$ are the Proj of the direct sum of these
by Morphisms, Lemma \ref{morphisms-lemma-proper-ample-is-proj}
we conclude that there is a canonical isomorphism $X_{n'} = Y$
as desired. We omit the verification that this is the unique
isomorphism making the diagram commute.
\end{proof}

\end{document}
